\scanpage{79}
Je désigne par $\widehat M^*_{\underline G,\underline G'}$, $\overset{\unsure{\deleted{\wedge}}}{M}{}^*_{\underline G,\underline G'}$ les germes de multiplicités obtenus en \deleted{enlevant (``l'infini'') $\widehat M_{\underline G}$ resp. $M_{\underline G}$, $\unsure{\Gamma^!_{\underline G,\underline G'}}$ (les traces de la \unsure{réunion} des diviseurs de $\unsure{\widehat M_{\underline G'}}$ \unsure{qui contiennent} $\unsure{(\widehat M_{\underline G},\Gamma^!_{G,G'})}$)}
\inserted{\deleted{\unsure{prenant les traces de} $M$} enlevant les traces des diviseurs à l'infini de $\widehat M_{\underline G'}$ \deleted{\unsure{non seuls}} qui contiennent $\widehat M_{\unsure{\underline G|G}}$ -- dans le cas de $M^*_{\underline G,\underline G'}$, on prend la \uline{trace} de $M_{\underline G'}$ sur $\widehat M_{\underline G,\underline G'}$}
Ils sont donc des filtres d'ouverts de la multiplicité $M_{G'}$, et les groupoïdes fondamentaux* \deleted{$\unsure{\Pi}$}
\[
(100)\qquad\Pi M^{\unsure{*}}_{\underline G,\underline G'},\qquad\Pi M^*_{\underline G,\underline G'}
\]
sont définis comme des Catégories $\underset{\longrightarrow}{\lim}$ des catégories de revêtements universels, \unsure{ou} ou comme des catégories des revêtts \unsure{du connexe} desdits germes de multiplicités (comparer [5], n°5, p. $768'$).

On a alors des diagrammes de germes de multiplicités analytiques
\diagram{\begin{tikzpicture}[x=1cm,y=1cm]
\node[inner sep=2pt] (A) at (0,0) {$M^*_{\underline G,\underline G'}$};
\node[inner sep=2pt] (B) at (4,0) {$M_{\underline G'}\deleted{\unsure{M}}$};
\node[inner sep=2pt] (C) at (-2.6,-2) {$(M_{\underline G}\deleted{\unsure{,\Gamma^!}})$};
\node[inner sep=2pt] (D) at (0,-2) {$M_{\underline G,\underline G'}$};
\draw[line cap=round] ([xshift=4pt]A.east) -- ([xshift=-8pt]B.west);
\draw[->] ([yshift=-5pt]A.south) -- ([yshift=5pt]D.north);
\draw[->] (C.east) -- (D.west);
\node[anchor=east] at (-3.75,-1.05) {$(101)$};
\node[anchor=north west,inner sep=0pt] (caption) at (-1.65,-2.55) {(équiv d'homotopie !)};
\node[anchor=base west,inner sep=0pt] at ([xshift=5pt,yshift=5pt]caption.north east) {$\unsure{)}$};
\end{tikzpicture}}
donnant lieu \unsure{à|:} des morphismes de groupoïdes
\margin{78}
